\section{Introduction}\label{sec:intro}

Let $X$ be a smooth projective complex threefold with $K_X\simeq\cO_X$
and $H^1(X,\cO_X)=H^2(X,\cO_X)=0$. For each effective class $\beta$, the characteristic function of the
compact space of cycles representing $\beta$ defines a class $U_{X,\beta}$ in
the Chern-number-zero, virtual-dimension-zero part $H$ of Pardon's rational
algebra of semi-Fano one-cycles \cite[\S\S2--3]{Pardon}. Write
$\Gamma_X\subset H_2(X;\ZZ)$ for the monoid of effective integral classes
and put $U_{X,0}=1$. We use the completion by degree against an ample
integral divisor. A bounded degree contains only finitely many effective
classes, by compactness of the Chow spaces and local constancy of the
integral cycle class.

There are two geometric descriptions of the local curves which enter
this problem. On a smooth curve $C$ of genus $g$, take a line bundle $L$
of degree $g-1$ and set $N=L\oplus(K_C\otimes L^{-1})$. The circle of
weights $+1,-1$ on these two summands fixes the multiple zero sections
$m[C]$. The associated equivariant Thom classes define elements
$a_{g,m}$ of $H$. Write $A_g(q)=1+\sum_{m\ge1}a_{g,m}q^m$.
The second description, due to Ionel and Parker, uses an almost-complex
structure for which every compact curve in a neighborhood of $C$ is
supported on $C$ and the normal operator is injective on every cover.
The contribution of such an elementary cluster is known numerically
\cite[\S3]{IP}. We compare the two descriptions on the space of cycles.

\begin{theorem}\label{thm:main}
For $X$ as above there are unique integers $n_{\beta,g}(X)$ indexed by
$\beta\in\Gamma_X\setminus\{0\}$ and $g\ge0$, zero for all
sufficiently large $g$ at each fixed $\beta$, such that
\begin{equation}\label{eq:product}
 \sum_{\beta\in\Gamma_X}U_{X,\beta}Q^\beta
 =\prod_{\beta\in\Gamma_X\setminus\{0\}}\prod_{g\ge0}
       A_g(Q^\beta)^{n_{\beta,g}(X)}
 \qquad\text{in }H\llbracket\Gamma_X\rrbracket .
\end{equation}
The integer $n_{\beta,g}(X)$ is the signed multiplicity of genus-$g$
elementary clusters of core class $\beta$ in a cluster decomposition
at any energy bound containing $\beta$. It equals the coefficient of
the genus-$g$ elementary series at $\beta$ in the connected
Gromov--Witten expansion of Ionel and Parker.
\end{theorem}

The balanced factors and their operations were constructed in
\cite[\S\S2--3]{PaperIII}. Theorem~1.1 of that paper obtained
\eqref{eq:product} assuming a universal Gromov--Witten evaluation on
analytic families. Here the integer exponents are determined by
deformation of cycles, before taking the connected numerical invariant.
The proof still uses compactness and gluing for stable maps of fixed genus.
Integral homology is used throughout, including when multiplication by
an integer identifies distinct classes.

The local comparison starts with a spin line $L$ and the real rotations
of the two copies of $L$ in $L\oplus L$. These rotations preserve the
Lee deformation of Ionel and Parker. On the integrable fiber the
circle has opposite weights. On each positive fiber the only cycle of
degree $m$ is $m[C]$. A compactly supported class on the total cycle
space restricts to the antidiagonal Thom class at one end and to the
normalized point class at the other (\S\ref{sec:local}).
One can see the effect of the choice of circle already at the double of
a genus-two curve. If $x_{g,m}$ denotes Pardon's scalar class, then
\[
                      a_{2,2}-x_{2,2}=-2x_{4,1}.
\]
In \S\ref{sec:double} we obtain the two negative terms separately from
the two families of branched double covers of a generic balanced curve.

To pass to $X$, we use the cluster isotopy of Ionel--Parker and
Doan--Ionel--Walpuski \cite[\S7]{IP}, \cite[\S3]{DIW}. In a bounded
energy range the disjoint connected supports can be separated in $X$.
A cycle with components in several of the resulting open sets gives a
product of compactly supported classes. At a fold two factors cancel
up to factors of higher multiples of the core class. Iteration gives a
finite signed product in an algebra of almost-complex families
(\S\S\ref{sec:families}--\ref{sec:clusters}).

The difficulty in returning to the complex cycle algebra is that a
relation among classes of curves can run over parameter spaces of
arbitrary dimension. Only the multiplicity-one classes need to be
distinguished. At the least genus in a putative relation the image-cycle
map from stable maps singles out one embedded curve, regular with
local degree one. Sections~\ref{sec:genus} and~\ref{sec:degree} show
that its point orientation survives in compactly supported cohomology.
A comparison of the first coefficient where two products differ then
gives \eqref{eq:product}. The result concerns the balanced antidiagonal
classes. The dependence of the other geometric local elements in
\cite[Remark~6.8]{Pardon} on their defining geometry is a separate question.

\subsection*{Acknowledgements}
I am grateful to John Pardon for a helpful discussion of the local curve classes.
